% page 384 du fac-simile (image p-383 du scan)
% bloc 152.4 x 237.7 mm ; marge gauche 8.0 mm
\pgc{-12.700mm}{0.000mm}{
\l{\cel{3}376}
\l{}
\l{miriapodo. (\sou{zool.)}\cel{2}Animalo artikoza (artropodo) qua kon-}
\l{\cel{3}sistas ye ringi, singla de qui havas paro de gambi. -}
\l{\cel{3}L.\cel{1}\sou{myriapoda}. - DEFIS.}
\l{}
\l{mirmekofago.\cel{2}(zool.) Mamifero ek la klaso \textquotedbl{}sen-denti\textquotedbl{},}
\l{\cel{3}di qua la lango esas filo-forma, viskoza, e quan lu ex-}
\l{\cel{3}tensas sur la formikeyo por kaptar la formiki qui esas}
\l{\cel{3}lua disho precipua. - L.\cel{1}\sou{myrmedophaga}. - IL.}
\l{}
\l{\sou{mir}mekofila. (\sou{bot.)}. Dicesas pri ula animali qui esas la}
\l{\cel{3}gasti di la formiki - qui vivas kun ici, sive kom \textquotedbl{}in-}
\l{\cel{3}vititi\textquotedbl{} (ex.: la afidii o planto-lausi), sive kom \textquotedbl{}ne-}
\l{\cel{3}atenciti\textquotedbl{}, sive kom vera paraziti. - DEFIS.}
\l{}
\l{\sou{mirmekoleono}. (\sou{zool.)}\cel{2}Insekto, analoga a libelulo, di qua}
\l{\cel{3}la larvo jacas sur la fundo di funelo quan lu exkavis}
\l{\cel{3}en tero e kaptas por manjar li la formiki od altra in-}
\l{\cel{3}sekti qui falas aden lu. - L.\cel{1}\sou{myrmecoleon}.-IL.}
\l{}
\l{\sou{mirto}. (\sou{bot.}) Arboreto kun flori blanka, quan la antiqui}
\l{\cel{3}konsakrabis a Venus (Afrodito). - DEFIRS. L.\cel{1}\sou{myrtus}.}
\l{}
\l{\sou{mirte}lo. (\sou{bot.)}\cel{3}Genero de planti ek \textquotedbl{}vacinei\textquotedbl{} qua pro-}
\l{\cel{3}duktas beri nigratra, acideta, akreta. - L.\cel{1}\sou{vaccinium}}
\l{\cel{3}\sou{murtidus}. - FIL.}
\l{}
\l{\sou{mis-}\cel{3}Prefixo qua signifikas \textquotedbl{}erore, nejuste\textquotedbl{}.}
\l{}
\l{\sou{misiono.}\cel{2}Senditaro (religiala, ciencala, diplomacala) di}
\l{\cel{3}qua la tasko konsistas ye - rispektive - predikar, ins-}
\l{\cel{3}traktar o docar, studiar, explorar, negociar, e c.)}
\l{\cel{3}- DEFIRS.}
\l{}
\l{\sou{mispelo.}\cel{1}(\sou{bot.}) Frukto qua havas plura kerni, qua manjesas}
\l{\cel{3}erste kande lu oldeskis e komencas velkar. - L.\cel{1}\sou{mes}-}
\l{\cel{3}\sou{pilus}. - DISL.}
\l{}
\l{\sou{mistelo}. (\sou{bot.)}\cel{2}Planto lignoza, qua vivas parazite sur}
\l{\cel{3}ula arbori (querko, pomiero, fraxino.) L.\cel{1}\sou{viscum.}}
\l{\cel{3}DE.}
\l{}
\l{\sou{misterio}. I. Rito sekreta di la politeismo antiqua, a qua}
\l{\cel{3}admisesis nur poka iniciiti. - To quon onu konservas,}
\l{\cel{3}mantenas, sekreta. - II. Dogmato kristana, revelita kom}
\l{\cel{3}objekto di kredo, e qua preter-iras raciono. - DEFIRS.}
\l{}
\l{\sou{mistifikar.}\cel{1}(trans.) Trompar, erorigar (ulu), amuzante su}
\l{\cel{3}de lua kredo. - DEFIR.}
\l{}
\l{\sou{mistiko.}\cel{1}Doktrino filozofiala, religiala, qua konsideras kom}
\l{\cel{3}la perfekteso sorto di kontemplo ed extazo qua elevas la}
\l{\cel{4}homo, ja en la vivo sur-tera, ad uniono misterioza kun}
\l{\cel{3}deo. - DEFIRS.}
}
